\documentclass[11pt,a4paper]{article}
\usepackage[margin=26mm]{geometry}
\usepackage{amsmath}
\usepackage{fontspec}
\setmainfont{Nimbus Roman}
\setsansfont{Nimbus Sans}
\setmonofont{Nimbus Mono PS}
\usepackage{mtpro2}
\usepackage{booktabs,tabularx,longtable,enumitem}
\usepackage[hidelinks,pdftitle={mtp2otf: Conversion and validation specification},pdfauthor={}]{hyperref}
\setlength{\parindent}{0pt}
\setlength{\parskip}{5pt}
\setlist{nosep,leftmargin=*}
\clubpenalty=10000
\widowpenalty=10000
\newcommand{\file}[1]{\nolinkurl{#1}}
\title{mtp2otf\\\large Conversion and validation specification}
\author{}
\date{}
\begin{document}
\maketitle
\vspace{-2em}

\section{Scope}
mtp2otf converts locally supplied MathTime\texttrademark{} Professional 2 (MTPro2) Type~1
fonts and TeX metrics into a single-face OpenType MATH font. The Full and
Lite editions have distinct installation identities. CFF-based OTF is the
canonical generated font; TrueType-based TTF is derived from it for Word.
Both formats retain the OpenType math layout data.

The interface to TeX engines is \file{unicode-math}. There is no companion
package, Lua callback, or runtime-generated TeX sidecar. The converter does
not resolve TeX trees, download inputs, install dependencies, or install
generated fonts. Full is the default. Missing Full inputs never select
Lite automatically.

\section{Inputs and build interface}
\subsection{Source staging}
Place the source files directly in \file{mtpro2/}. Each active source font
requires its PFB and TFM pair; \file{mtpro2.sty} supplies package policy.
Full additionally requires \file{umt2ms.fd}, \file{umt2mf.fd},
\file{umt2bb.fd}, and \file{umt2hrb.fd}. Lite does not require these
Full-only declarations. The source registry and required-file calculation
are defined in \file{tools/source_policy.py}. Empty inputs receive setup
guidance; partially supplied inputs receive a missing-file list.

The Lite source set contains the text, script, and scriptscript sizes of
\file{mt2mi}, \file{mt2sy}, and \file{mt2mb}, plus
\file{mt2exa}, \file{mt2exe}, \file{mt2exf}, \file{mt2exg},
\file{mt2xl}, and \file{mt2xxxl}. The three optical-size suffixes are
\file{t}, \file{s}, and \file{f}. Full includes the additional math
alphabets and symbol sources in the registry.

\subsection{Roman donors}
The default inputs in \file{donors/} are
\file{NimbusRoman-Regular.otf} and \file{NimbusRoman-Bold.otf}
from Artifex urw-base35. Another static CFF OTF or TrueType TTF pair can
be selected explicitly. Both faces must be upright, proportional, from
the same family, and have the required Latin and digit coverage and
appropriate weights. Coordinates and advances are normalized to the
output units per em.

Regular supplies upright Latin and permitted additional text characters;
Bold supplies mathematical bold Latin and digits. Mathematical italic,
bold italic, and non-Latin-shaped Greek come from MTPro2. Optional donor
text coverage can vary. A donor must not fill missing Full-only math
semantics in Lite.

\subsection{Commands and configuration}
\begin{verbatim}
mkdir -p mtpro2 donors
# Copy the source fonts, declarations, and donors.
./build.sh
./build.sh --edition=lite
./build.sh --roman-regular=/path/to/Regular.otf \
           --roman-bold=/path/to/Bold.otf
\end{verbatim}
The public options are \file{--edition}, \file{--roman-regular}, and
\file{--roman-bold}; help is available with \file{--help}.
Unknown options and abbreviated option names are rejected.
Relative donor paths are resolved from the caller's working directory.
All project directories are resolved from the location of \file{build.sh}.

Python~3.10 or later with fontTools, FontForge with Python~3.10 or later
support, and \file{tftopl} are required. Dependency and input failures stop
the build early. Users create and populate the input directories; the build
creates only its working and output directories.

\section{Working state and final outputs}
\begin{tabularx}{\linewidth}{@{}l>{\raggedright\arraybackslash}X@{}}
\toprule
Location & Purpose \\
\midrule
\file{mtpro2/} & Local Type~1 fonts, TeX metrics, and declarations. \\
\file{donors/} & Local Regular and Bold Roman donors. \\
\file{build/full/}, \file{build/lite/} & Edition-specific intermediate
fonts, converted PL files, calculated values, and one diagnostic log. \\
\file{out/} & Validated OTF and TTF deliverables only. \\
\bottomrule
\end{tabularx}

Full produces \file{MTPro2Math.otf} and \file{MTPro2Math.ttf}.
Lite produces \file{MTPro2MathLite.otf} and \file{MTPro2MathLite.ttf}.
Both editions coexist in \file{out/}.
No font is built directly into this directory.

After every required check succeeds, each validated candidate is copied
to a temporary file beside its destination and renamed into place.
This avoids exposing a partially copied individual font. It is not a
two-file filesystem transaction: an interruption during final promotion
can leave a previously validated OTF/TTF pair from different builds.
Run the build again to replace the pair together in the normal sequence.

Input, construction, conversion, or validation failure leaves previous
outputs intact. Their presence does not establish that the latest build
succeeded: use the command's exit status. Full and Lite may be built in
sequence; simultaneous builds of the same edition are not supported.
The entire \file{build/} directory is disposable and can be regenerated
from the local inputs.

\subsection{Progress and diagnostics}
Normal stdout contains one line: \texttt{building MTPro2 Math}, or
\texttt{building MTPro2 Math Lite}. Successful audits are silent. An unsuccessful
subprocess reports its captured diagnostics on stderr. FontForge output
and warnings remain available in \file{build/<edition>/build.log} even
when the subprocess succeeds. The log is replaced for each new build
that reaches this stage, rather than accumulated across runs.

Exit status is zero on success, two for command usage and preflight
failures handled by the frontend, and one for build or validation failure.
There is no automatic installation, clean command, verification bypass,
or verbose-mode configuration.

\section{Conversion contracts}
\subsection{Source metrics and ordinary spacing}
An untransformed MTPro2 glyph retains TFM \file{CHARWD} as its horizontal
advance and \file{CHARIC} as its MATH italic correction. The ordinary
adjacency adjustment is
\[
  \operatorname{CHARIC}(\text{left})
  + \operatorname{TFMkern}(\text{left},\text{right}).
\]
It is encoded in cumulative GPOS layers. A missing TFM pair contributes
zero, not a reason to discard italic correction. Skewchar-right pairs
are excluded from ordinary kerning according to the local declarations.
Script, prime, and accent placement are separate layout paths.

For source math alphabets, top-accent attachment is based on
\[
  \frac{\operatorname{CHARWD}+\operatorname{CHARIC}}{2}
  + \operatorname{skewkern}.
\]
Combining accents and synthesized forms follow their own coordinate
transformations. Widths and attachment coordinates are not interchangeable.

\subsection{Optical sizes, aliases, and variants}
Text, script, and scriptscript source outlines are exposed through
OpenType math substitutions. Full-only option alphabets preserve their
optical-size reachability in both supported feature application orders.
Latin-shaped Greek aliases share glyph IDs with their canonical donors;
the aliases do not duplicate outlines or acquire separate metrics.

TFM larger-size chains, extensible recipes, and local outline geometry
determine variant and assembly construction. Repeatable parts retain
both connector endpoints. Brace cusp geometry follows the source halves;
the cusp inside a part is distinct from an overlap between assembly parts.

\subsection{Primes and compatibility forms}
The public prime forms and contextual substitutions provide the retained
raw-prime and superscript behavior. Their padding is derived from local
outline geometry. Scriptscript prime alternates use the native optical
source forms, with their corresponding source widths. Backprime follows
its own source identity. These are conversion policies, not a claim that
every prime transformation is a direct TFM assignment.

Public vector and bar accents are fixed. Private wide forms are distinct:
U+E286 supplies finite variants, while U+E287 supplies an assembly.
Standard commands do not automatically select these private paths.
Swash z remains available through \file{cv01} or U+E000 where supported;
there is no default swash or Word command registration.

\subsection{MATH constants and provenance}
\file{tools/math_constants.py} calculates the MATH constants from local
source metrics and explicitly retained conversion policies.
\file{tools/source_policy.py} extracts package and font-definition policy
and the geometry needed by synthesized forms. The calculation covers
the 56 MathConstants fields and the separate common connector overlap.

Not every field has a direct TFM counterpart. Reference-derived
thresholds and engine-compatibility choices are documented as such;
they are not relabelled as original-source metrics. In particular,
upper-limit placement retains a compatibility policy rather than a
universal one-to-one TeX parameter mapping. Next-style baseline-drop
values are converted into the parent font's coordinate scale.

The appendix gives every scalar assignment and its source classification.
It describes formulas and explicit conversion policy, not evaluated local
MTPro2 measurements. The same assignments apply to Full and Lite.

The symbol TFM supplies axis height and script/fraction shifts; the
extension TFM supplies rule thickness and large-operator spacing. The
package supplies script ratios. TeX-style gap formulas combine those
inputs. OpenType-only thresholds and skewed-fraction gaps are retained
migration policies rather than recovered TFM parameters. Rendered
comparisons support their retention in tested engines, not universal
identity with every original TeX expression.

Ordinary superscripts select the source sup2 branch and subscripts select
sub2; OpenType's common fields do not encode every TeX style branch.
Upper-limit and stretch-stack upper-side assignments preserve the reviewed
Word placement policy. Lower stretch-stack controls reuse the lower-limit
inputs and are not claimed to be independently optimized. Baseline drops
use script-size symbol parameters converted to the first-level parent's
em; nested engine scaling can differ. The flattened-accent threshold does
not imply a dedicated flattened-accent substitution.

\section{Validation methodology}
Validation is mandatory. Reuse of parsing and identity definitions does
not replace comparisons of serialized output with local source metrics.
OTF establishes source semantics. Cross-format invariants transfer those
guarantees to TTF; structural and naming checks run on both formats.

\begin{tabularx}{\linewidth}{@{}l>{\raggedright\arraybackslash}X@{}}
\toprule
Check & Responsibility \\
\midrule
Source declarations & Required source set, unique usable policy,
accent semantics, and Full optical-family ordering. \\
Roman donors & Input type, family, weight, required coverage; output
bounding boxes and advances, plus an extra-import heuristic. \\
Structural validator & Tables, cmap, aliases, clipping bounds,
assemblies, features, prime contexts, and serialization. \\
Edition naming & Installation names, technical version metadata,
and CFF naming constraints for both formats. \\
Source contract & Advances, italic correction, top accents, and
ordinary IC plus TFM pair adjustments in the OTF; converted accent
centers are checked against TTF outlines. \\
Edition capabilities & Full option-alphabet optical sizes or
Lite's required and forbidden features and code points. \\
Cross-format check & Glyph-ID mapping, advances and sidebearings,
MATH/GSUB/GPOS/GDEF table bytes, and outline origins. \\
MATH values & Each calculated scalar against the serialized OTF,
including the common connector overlap. \\
\bottomrule
\end{tabularx}

The Roman output check compares normalized bounding boxes and advances;
it is not a proof of contour identity. The OTF/TTF check permits the
intended cubic-to-quadratic approximation and checks origin shifts
separately. Engine rendering is not proven solely by table validation.

\subsection{Local evidence}
\file{build/<edition>/source-policy.json} contains parsed policy,
geometry, and input hashes. \file{mtpro2-source-contract.json} contains
the glyph relationships and source-derived records used by output
audits. \file{math-constants.json} contains evaluated scalars and their
provenance descriptions. These snapshots are working data, not extra
deliverables. Repeated PASS summaries and copies of the same snapshots
are not written beside the final fonts.

When changing the converter, build both editions from the same inputs
before and after the change. Compare corresponding OTFs and TTFs with
build timestamps and dependent checksums normalized; retain all other
table differences for investigation. A layout change additionally
requires focused engine tests. Identical nonvolatile output can reuse
the earlier rendering evidence for that exact font content.

\section{Specimens and engine limits}
\file{word.docx} and \file{lualatex.tex} are editable specimens. The
included PDFs show the Full edition with Nimbus donors. Word uses the
TTF as its math font and Times New Roman for body text; LuaLaTeX uses
the OTF and Nimbus Roman for body text. Close Word before replacing
an installed font.

To rebuild the LuaLaTeX specimen, first build Full, then run
\texttt{lualatex lualatex.tex} from \file{doc/}. The specimen needs
\file{amsmath}, \file{fontspec}, \file{unicode-math}, and all four Nimbus
Roman body styles discoverable by TeX. These specimen dependencies are
not additional converter dependencies. This specification itself can
be compiled with \texttt{pdflatex mtp2otf.tex}, without MTPro2 fonts.

Word can segment mathematical runs before GPOS processing. LuaHBTeX
may not reproduce the complete Type~1 ordinary italic-correction plus
pair-kern combination. A correct font table therefore does not imply
identical spacing in every engine. No Lua option or callback is added
to hide these differences. Compare the included specimens with the original
Type~1 output in the same layout context when distinguishing source
behavior, conversion behavior, and engine behavior.

\section{Publication and references}
The public tree contains the shell entry point, Python tools, README,
LICENSE, and documentation. Local input and working directories are
not part of the archive. Publication inspection is available as
\texttt{python3 tools/audit\_public\_source.py .} on a clean source-only
tree. It checks prohibited paths and selected metric patterns; it is
not a general legal or confidentiality audit of arbitrary PDF contents.

Mappings are interoperability metadata created for this project by
comparing TeX and Unicode semantics, local-source behavior, and rendered
results. Type~1 names alone do not determine the required mappings.
They are not mechanical copies of MTPro2 files.

Project-authored code and documentation use 0BSD. MTPro2, donor fonts,
and external tools retain their own terms. MathTime is a trademark of
Publish or Perish, Inc. This is an independent project, not affiliated
with or endorsed by Publish or Perish, Inc., Personal TeX, Inc., or
Artifex Software.

\begin{itemize}
\item Repository and implementation:
\url{https://github.com/satom-zk/mtp2otf}
\item Bug reports:
\url{https://github.com/satom-zk/mtp2otf/issues}
\end{itemize}

\appendix
\section{Scalar assignments}
The following expressions follow \file{tools/math_constants.py} in its
actual evaluation order. \texttt{round} is Python nearest-integer rounding
(ties to even); \texttt{//} is floor division. \texttt{Fraction(a,b)} is an
exact rational. Intermediate rounding is intentional and must not be
algebraically reordered.

\texttt{upm} is the output units per em. \texttt{s(i)} means
\texttt{round(sy[i] * upm)}, with \texttt{sy} from \file{mt2syt.tfm};
\texttt{e(i)} uses \file{mt2exa.tfm} in the same way.
\texttt{theta = e(8)} and \texttt{xh = s(5)}.
\texttt{script\_sy} comes from \file{mt2sys.tfm}.
The two ratios are the local \file{mtpro2.sty} defaults;
\texttt{\_pct(r)} evaluates
\texttt{round(r.numerator * 100 / r.denominator)}.

\texttt{display\_operator\_min\_height} is the floor of the average of
the separately rounded height-plus-depth values of extension TFM slots
\texttt{0x50} and \texttt{0x58}, in output units. The source provides the
sizes; choosing their midpoint is a conversion policy. The delimited
subformula threshold reuses this value and is not TeX's \texttt{delim1}.

Source classes distinguish inputs from assignment choices:
MT-TFM denotes local TFM parameters; MT-STY denotes package declarations;
MT-PFB denotes local outlines; OT-SUG denotes an OpenType suggested
formula; TeX denotes the script-space convention. OT-MIG and MT-MIG denote
explicit migration policy. Combined labels mean that source data and an
assignment policy both contribute. A source label does not establish
exact layout equivalence in every engine.

Values are in font units except the two script percentages and the
radical-degree percentage. Numeric literals in migration rows are project
policy, not a table of measurements copied from MTPro2. The initial
connector-overlap value is reduced to the nonnegative integer minimum
joining run when local geometry requires a smaller overlap. Both ends of
repeatable parts participate; fixed outside ends do not.

\begingroup
\small
\setlength{\tabcolsep}{4pt}
\begin{longtable}{@{}p{0.46\linewidth}p{0.48\linewidth}@{}}
\toprule
Field & Expression and source class \\
\midrule
\endfirsthead
\toprule
Field & Expression and source class \\
\midrule
\endhead
\bottomrule
\endfoot
\texttt{Axis\allowbreak Height} & \texttt{s(22)}\newline MT-TFM \\
\texttt{Fraction\allowbreak Numerator\allowbreak Shift\allowbreak Up} & \texttt{s(9)}\newline MT-TFM \\
\texttt{Fraction\allowbreak Numerator\allowbreak Display\allowbreak Style\allowbreak Shift\allowbreak Up} & \texttt{s(8)}\newline MT-TFM \\
\texttt{Fraction\allowbreak Denominator\allowbreak Shift\allowbreak Down} & \texttt{s(12)}\newline MT-TFM \\
\texttt{Fraction\allowbreak Denominator\allowbreak Display\allowbreak Style\allowbreak Shift\allowbreak Down} & \texttt{s(11)}\newline MT-TFM \\
\texttt{Fraction\allowbreak Numerator\allowbreak Gap\allowbreak Min} & \texttt{theta}\newline OT-SUG+MT \\
\texttt{Fraction\allowbreak Num\allowbreak Display\allowbreak Style\allowbreak Gap\allowbreak Min} & \texttt{3 \allowbreak * theta}\newline OT-SUG+MT \\
\texttt{Fraction\allowbreak Denominator\allowbreak Gap\allowbreak Min} & \texttt{theta}\newline OT-SUG+MT \\
\texttt{Fraction\allowbreak Denom\allowbreak Display\allowbreak Style\allowbreak Gap\allowbreak Min} & \texttt{3 \allowbreak * theta}\newline OT-SUG+MT \\
\texttt{Fraction\allowbreak Rule\allowbreak Thickness} & \texttt{theta}\newline MT-TFM \\
\texttt{Stack\allowbreak Top\allowbreak Shift\allowbreak Up} & \texttt{s(10)}\newline MT-TFM \\
\texttt{Stack\allowbreak Top\allowbreak Display\allowbreak Style\allowbreak Shift\allowbreak Up} & \texttt{s(8)}\newline MT-TFM \\
\texttt{Stack\allowbreak Bottom\allowbreak Shift\allowbreak Down} & \texttt{s(12)}\newline MT-TFM \\
\texttt{Stack\allowbreak Bottom\allowbreak Display\allowbreak Style\allowbreak Shift\allowbreak Down} & \texttt{s(11)}\newline MT-TFM \\
\texttt{Stack\allowbreak Gap\allowbreak Min} & \texttt{3 \allowbreak * theta}\newline OT-SUG+MT \\
\texttt{Stack\allowbreak Display\allowbreak Style\allowbreak Gap\allowbreak Min} & \texttt{7 \allowbreak * theta}\newline OT-SUG+MT \\
\texttt{Superscript\allowbreak Shift\allowbreak Up} & \texttt{s(14)}\newline MT-TFM+OT-MIG \\
\texttt{Superscript\allowbreak Shift\allowbreak Up\allowbreak Cramped} & \texttt{s(15)}\newline MT-TFM \\
\texttt{Subscript\allowbreak Shift\allowbreak Down} & \texttt{s(17)}\newline MT-TFM+OT-MIG \\
\texttt{Subscript\allowbreak Top\allowbreak Max} & \texttt{round(Fraction(4,\allowbreak 5) \allowbreak * s(5))}\newline OT-SUG+MT \\
\texttt{Superscript\allowbreak Bottom\allowbreak Min} & \texttt{round(Fraction(1,\allowbreak 4) \allowbreak * s(5))}\newline OT-SUG+MT \\
\texttt{Superscript\allowbreak Bottom\allowbreak Max\allowbreak With\allowbreak Subscript} & \texttt{round(Fraction(4,\allowbreak 5) \allowbreak * s(5))}\newline OT-SUG+MT \\
\texttt{Sub\allowbreak Superscript\allowbreak Gap\allowbreak Min} & \texttt{4 \allowbreak * theta}\newline OT-SUG+MT \\
\texttt{Upper\allowbreak Limit\allowbreak Baseline\allowbreak Rise\allowbreak Min} & \texttt{e(9)}\newline OT-MIG+MT \\
\texttt{Upper\allowbreak Limit\allowbreak Gap\allowbreak Min} & \texttt{e(11)}\newline OT-MIG+MT \\
\texttt{Lower\allowbreak Limit\allowbreak Gap\allowbreak Min} & \texttt{e(10)}\newline MT-TFM \\
\texttt{Lower\allowbreak Limit\allowbreak Baseline\allowbreak Drop\allowbreak Min} & \texttt{e(12)}\newline MT-TFM \\
\texttt{Stretch\allowbreak Stack\allowbreak Top\allowbreak Shift\allowbreak Up} & \texttt{e(9)}\newline OT-MIG+MT \\
\texttt{Stretch\allowbreak Stack\allowbreak Gap\allowbreak Above\allowbreak Min} & \texttt{e(11)}\newline OT-MIG+MT \\
\texttt{Stretch\allowbreak Stack\allowbreak Gap\allowbreak Below\allowbreak Min} & \texttt{e(10)}\newline OT-MIG+MT \\
\texttt{Stretch\allowbreak Stack\allowbreak Bottom\allowbreak Shift\allowbreak Down} & \texttt{e(12)}\newline OT-MIG+MT \\
\texttt{Overbar\allowbreak Rule\allowbreak Thickness} & \texttt{theta}\newline MT-TFM \\
\texttt{Overbar\allowbreak Vertical\allowbreak Gap} & \texttt{3 \allowbreak * theta}\newline OT-SUG+MT \\
\texttt{Overbar\allowbreak Extra\allowbreak Ascender} & \texttt{theta}\newline OT-SUG+MT \\
\texttt{Underbar\allowbreak Rule\allowbreak Thickness} & \texttt{theta}\newline MT-TFM \\
\texttt{Underbar\allowbreak Vertical\allowbreak Gap} & \texttt{3 \allowbreak * theta}\newline OT-SUG+MT \\
\texttt{Underbar\allowbreak Extra\allowbreak Descender} & \texttt{theta}\newline OT-SUG+MT \\
\texttt{Radical\allowbreak Rule\allowbreak Thickness} & \texttt{theta}\newline MT-TFM \\
\texttt{Radical\allowbreak Extra\allowbreak Ascender} & \texttt{theta}\newline OT-SUG+MT \\
\texttt{Radical\allowbreak Vertical\allowbreak Gap} & \texttt{theta + theta // 4}\newline OT-SUG+MT \\
\texttt{Radical\allowbreak Display\allowbreak Style\allowbreak Vertical\allowbreak Gap} & \texttt{theta + xh // 4}\newline OT-SUG+MT \\
\texttt{Radical\allowbreak Kern\allowbreak Before\allowbreak Degree} & \texttt{round(Fraction(5,\allowbreak 18) \allowbreak * upm)}\newline OT-SUG \\
\texttt{Radical\allowbreak Kern\allowbreak After\allowbreak Degree} & \texttt{-round(Fraction(10,\allowbreak 18) \allowbreak * upm)}\newline OT-SUG \\
\texttt{Radical\allowbreak Degree\allowbreak Bottom\allowbreak Raise\allowbreak Percent} & \texttt{60}\newline OT-SUG \\
\texttt{Script\allowbreak Percent\allowbreak Scale\allowbreak Down} & \texttt{\_pct(script\_ratio)}\newline MT-STY \\
\texttt{Script\allowbreak Script\allowbreak Percent\allowbreak Scale\allowbreak Down} & \texttt{\_pct(scriptscript\_ratio)}\newline MT-STY \\
\texttt{Min\allowbreak Connector\allowbreak Overlap} & \texttt{100}\newline OT-MIG+MT-PFB \\
\texttt{Accent\allowbreak Base\allowbreak Height} & \texttt{477}\newline OT-MIG \\
\texttt{Flattened\allowbreak Accent\allowbreak Base\allowbreak Height} & \texttt{656}\newline OT-MIG \\
\texttt{Delimited\allowbreak Sub\allowbreak Formula\allowbreak Min\allowbreak Height} & \texttt{int(display\_operator\_min\_height)}\newline MT-MIG+MT \\
\texttt{Display\allowbreak Operator\allowbreak Min\allowbreak Height} & \texttt{int(display\_operator\_min\_height)}\newline MT-TFM \\
\texttt{Math\allowbreak Leading} & \texttt{154}\newline OT-MIG \\
\texttt{Skewed\allowbreak Fraction\allowbreak Horizontal\allowbreak Gap} & \texttt{350}\newline OT-MIG \\
\texttt{Skewed\allowbreak Fraction\allowbreak Vertical\allowbreak Gap} & \texttt{102}\newline OT-MIG \\
\texttt{Space\allowbreak After\allowbreak Script} & \texttt{round(Fraction(1,\allowbreak 20) \allowbreak * upm)}\newline TeX \\
\texttt{Subscript\allowbreak Baseline\allowbreak Drop\allowbreak Min} & \texttt{round(script\_sy[19] \allowbreak * float(script\_ratio) \allowbreak * upm)}\newline MT-TFM+OT-MIG \\
\texttt{Superscript\allowbreak Baseline\allowbreak Drop\allowbreak Max} & \texttt{round(script\_sy[18] \allowbreak * float(script\_ratio) \allowbreak * upm)}\newline MT-TFM+OT-MIG \\
\end{longtable}
\endgroup
\end{document}
